\begin{proof}
\begin{enumerate}[label=\arabic*.]
\item We first remove the unit matrix $C(q)$ topologically, by an explicit
deformation, without changing the integral marking.  Use the family
$C_u(q)=(1-u)C(q)+u iI_3$, followed by the reverse of the same
family starting at $C(q)=0$.  The two flows compose to a marked
comparison with $C(q)=0$.  The action and properness estimates
of P1.9 are uniform on both parameter intervals.  Local quotient charts form a smooth
family of pairs, with the local equation $q=x_0x_1x_2x_3$ independent
of $u$.  Local lifts of $\partial/\partial u$, tangent to every
coordinate stratum and annihilating $q$, can be patched by a partition
of unity.  Over a smaller compact disk their flow exists for the whole
interval.  This gives a stratified diffeomorphism of the central
fibres and a diffeomorphism of the nearby fibres, preserving the
period marking.  The same construction on the real oriented blowup
preserves the collapse map.  We may therefore calculate at $C(q)=0$.

\item Here is a concrete topological model.  Put $G=4I_3-\mathbf1\mathbf1^t$.
Its positive eigenvalues are $1,4,4$.  The Voronoi polytope in dual
coordinates is cut out by
\[
|v^ty|\leq \tfrac12v^tGv
 \quad (0\ne v\in\{0,1\}^3).
\]
The value on a binary vector with $k$ nonzero entries is $k(4-k)$.
These fourteen primitive facet normals give the star of a vertex in
the Kuhn triangulation.  One checks its Delaunay tetrahedra by the
successive chains
$0,e_i,e_i+e_j,e_1+e_2+e_3$ for the six permutations $(i,j,k)$;
solving their three perpendicular-bisector equations and the fourteen
inequalities verifies precisely their common faces.  The translates
of this Voronoi polytope by $G\mathbb Z^3$ tile $\mathbb R^3$.

\item For a simplex $\sigma$ with vertices $v_0,\ldots,v_r$, put
$L_\sigma=\sum_i\mathbb Z(v_i-v_0)$ and
$D_\sigma=\bigcap_{v\in\sigma}P_v$, where $P_v=P+Gv$.
The component torus has lattice $A=\mathbb Z^3$; the common-stratum
marking is the explicit map
\[
 (A\oplus\mathbb Z)/\sum_i\mathbb Z(v_i,1)
 \longrightarrow A/L_\sigma,\qquad
 [(a,k)]\longmapsto a-kv_0.
\]
The Kuhn unimodularity makes $L_\sigma$ primitive.

We choose compatible moment maps as follows.  The polytope $2P_v$
is integral, since its vertex facet normals are unimodular and its
facet constants are integral.  Embed the component $X_v$ using the
characters $\chi^{-m}$ for
$m\in\mathcal A_v=2P_v\cap A^\vee$, and define
\[
 \mu_v([Z_m])=
 \frac{\sum_{m\in\mathcal A_v}m|Z_m|^2}
 {2\sum_{m\in\mathcal A_v}|Z_m|^2}.
\]
The sign is specified: the ray $w-v$ tends to the face on which
$\langle w-v,-m\rangle$ is minimal, hence to the upper common
Voronoi face of $P_v$ and $P_w$.  Thus the character polytope is
$-2P_v$, and $\mu_v$ is minus one-half its usual moment map.
The monomial moment theorem
\cite[equations (12.2.3)--(12.2.4) and Theorem 12.2.5, p.~571]{CoxLittleSchenck2011}
identifies the nonnegative toric part with $P_v$.
For the squared weights here, the intervening map on that part is
$[a_m]\mapsto[a_m^2]$, with inverse $[b_m]\mapsto[\sqrt{b_m}]$:
both preserve the homogeneous binomial relations and positive common
projective rescaling.  Write $s_v:P_v\to X_v$ for the resulting
continuous positive section.

On a shared face put
$\mathcal A_\sigma=2D_\sigma\cap A^\vee$ and choose
$m_0\in\mathcal A_\sigma$.
For $m\in\mathcal A_\sigma$, the difference $m-m_0$ annihilates
$L_\sigma$.  The ratio $Z_m/Z_{m_0}$ is consequently the restriction
of the same character
\[
 \bigl(-(m-m_0),\langle m-m_0,v\rangle\bigr)
       \in(A\oplus\mathbb Z)^\vee
\]
for every vertex $v$ of $\sigma$.  At $C=0$ all deck unit multipliers are one, and the toric gluing
identifies these character functions exactly.  On either component
the restricted moment map is therefore the identical expression
\[
 \frac{\sum_{m\in\mathcal A_\sigma}
       m|\chi^{-(m-m_0)}|^2}
 {2\sum_{m\in\mathcal A_\sigma}|\chi^{-(m-m_0)}|^2}.
\]
The equality persists on every smaller face by restricting the index
set.  Positive character values also agree, so the sections $s_v$
glue on the entire shared orbit closure.  The angular maps are the
displayed quotient homomorphisms $A\to A/L_\sigma$.

The period deck map is $(a,k)\mapsto(a+k\nu,k)$, sending $v$ to
$v+\nu$.  Its component marking is unchanged because
$a+k\nu-k(v+\nu)=a-kv$.
Moreover $\mathcal A_{v+\nu}=\mathcal A_v+2G\nu$; the common
character factor cancels in projective ratios.  Hence
\[
 \mu_{v+\nu}(T_\nu p)=\mu_v(p)+G\nu,
 \qquad s_{v+\nu}(y+G\nu)=T_\nu s_v(y).
\]
These give an explicit well-defined map
\[
 \begin{gathered}
 \mathcal Q=
 \bigl((\mathbb R^3/G\mathbb Z^3)\times
       (\mathbb R^3/\mathbb Z^3)\bigr)/\!\sim
 \longrightarrow C,\\
 [y,\theta]\longmapsto e^{2\pi i\theta}\,s_v(y)
 \quad(y\in P_v).
 \end{gathered}
\]
where over $\operatorname{relint}D_\sigma$ the relation collapses
the subtorus $L_\sigma\otimes(\mathbb R/\mathbb Z)$.
On each stratum the inverse takes its moment value and its angular
character class in $T_A/T_{L_\sigma}$; the latter is obtained by
$z\mapsto z/|z|$ on its nonzero characters.
The map is a continuous bijection, and its source is compact and
$C$ Hausdorff, so it is a homeomorphism.
The angular identity $[y,0]$ gives the continuous section, with
the same face and deck compatibilities just proved.

\item The actual nearby-fibre collapse is the quotient map in this model.
Give the semistable family its divisorial logarithmic structure.
The local equations $q=x_0x_1x_2x_3$ and their face restrictions
make it fine, saturated, log smooth and vertical over the standard
logarithmic disc.  The family is proper by P1.9.
In its logarithmic central fibre, retain all arguments of the $x_j$
with their sum fixed by $\arg q$.  The integral angular lattice over
a face with vertices $v_0,\ldots,v_r$ is the kernel of the height
map on $\mathbb Z^3\oplus\mathbb Z$.  Forgetting the arguments of
the zero coordinates induces exactly
\[
 \mathbb Z^3\longrightarrow
 (\mathbb Z^3\oplus\mathbb Z)/
       \sum_i\mathbb Z(v_i,1)
 \ \cong\ \mathbb Z^3/L_\sigma.
\]
The last isomorphism sends the class of $(a,k)$ to
$a-kv_0\pmod {L_\sigma}$; its inverse sends $a$ to $(a,0)$.
Thus no finite index enters this angular quotient.
The polar maps agree under the specified integral monomial overlaps.
At $C(q)=0$ the deck multipliers have positive real values, so the
logarithmic fibre at angle one is the product of the base and angular
three-tori in item~3, with their given markings.

The mapping-cylinder comparison for a proper vertical log-smooth
family identifies this logarithmic forgetful map with the actual
nearby-fibre collapse, compatibly with the central strata
\cite[formula (3-3-1), p.~478; Theorem 4.1.1 and Lemmas
4.1.2--4.1.3, pp.~480--482]{ffAchingerOgus20}.
Its angular cycles are $a_j$ and its base cycles are $b_j$.
In these column-basis conventions monodromy fixes each $a_j$
and sends $b_j$ to $b_j+a_j$, as specified by $Nb_j=a_j$.
The stratified deformation in item~1 transports this actual map
back to the original bounded period matrix.

\item Filter $\mathcal Q$ by the dimension $d$ of its dual base cells.
Over the interior of the dual cell of a $p$-simplex, $d=3-p$ and the
angular quotient has lattice
$Q_\sigma=\mathbb Z^3/L_\sigma$ of rank $d$.
Excision for the cell and its boundary gives the relative chain
groups of the resulting homology spectral sequence:
\[
E^1_{d,q}=\bigoplus_{\dim\sigma=3-d}\bigwedge^qQ_\sigma.
\]
Orient dual cells by the orientation of the base and the ordered
simplex.  The differential is the signed quotient map for each
codimension-one coface.  More explicitly, with integral projection
matrices $Q_\sigma$ and chosen right inverses $S_\sigma$,
\[
D_{p,q}[\tau,\sigma]=
 \sum_{\sigma=\partial_j\tau\ {\rm mod}\ \mathbb Z^3}
 (-1)^j\bigwedge^q(Q_\tau S_\sigma).
\tag{1.11.1}
\]
All $L_\sigma$ are primitive because the tetrahedra are unimodular.
Thus these projections and right inverses are integral.  Formula
(1.11.1) specifies every integer entry, including multiple incidences
under translation identifications.  The program
\texttt{cusp\_complex.py} enumerates the $26$ simplex classes, all
$74$ codimension-one incidences, and exports every matrix, basis, and
Smith certificate to \texttt{cusp\_complex.json}.

\item The nonzero $E^2_{d,q}$ ranks are
\[
\begin{array}{c|rrrr}
q\backslash d&0&1&2&3\\\hline
0&1&3&3&1\\
1&0&6&8&3\\
2&0&0&7&3\\
3&0&0&0&1
\end{array}
\]
and every group is free.  Support leaves only the possible higher
differential $d_2:E^2_{3,0}\to E^2_{1,1}$.
The angular identity section splits the base projection as a map of
filtered spaces, so the base fundamental class survives.  This kills
that differential.  All subsequent differentials have zero source or
target.  The filtration extensions are extensions of free abelian
groups and therefore split, proving the asserted integral homology.

\item The product-to-quotient map induces on $E^1$ the exterior powers of
the quotient projections.  The newly authored program
\path{P1_11_collapse_fresh.py} in \path{code/continuation_V/}
computes the induced maps on integral cycle and quotient bases.
Their images are primitive, and the only nonzero graded cokernel
is $\mathbb Z$ at $(d,q)=(2,2)$.  The summed image ranks are
\[
             r_k=(1,3,9,9,9,3,1).
\]
The same exact exterior-power calculation gives these ranks for
$\ker(\bigwedge^kM_0-I)$.
The actual collapse satisfies
$\kappa_*(\bigwedge^kM_0-I)=0$, because a monodromy circuit
has the same retraction to the central fibre in the mapping-cylinder
comparison.  Therefore its rank is at most $r_k$.
A filtered map has rank at least the sum of the ranks of its
graded maps; hence the actual rank is exactly $r_k$.

Choose integral bases compatible with the source and target
filtrations; their graded groups are free.  In each graded block
choose the unit Smith pivots of the computed map.  The corresponding
$r_k$ rows and columns of the actual filtered map form a block
triangular submatrix with unimodular diagonal pivot blocks.
Its determinant is $\pm1$.  Since the full map has rank $r_k$,
these unit pivots prove that its actual image is primitive.
Thus the actual map is onto in every degree except four, where its
cokernel is the free group of rank $10-9=1$.
This rank argument also rules out additional lower-filtration images
of source classes in the kernels of the graded maps.

\item For a fully marked certificate we represent a base cycle $b_J$ by
Poincar\'e duality from the Alexander--Whitney cup cocycle on the
complementary coordinates.  On an ordered $p$-simplex its coefficient
is
\[
\operatorname{sgn}(J,J^c)
 \prod_{r=1}^p(v_r-v_{r-1})_{(J^c)_r}.
\]
Tensoring with an angular exterior vector and applying (1.11.1)
gives the matrices exported by \texttt{cusp\_marked\_collapse.py}.
They annihilate $\bigwedge^kM_0-I$ exactly.  The actual collapse also
annihilates that operator, since turning the smoothing parameter
does not change its image in the central fibre.  The computed ranks
equal the ranks of the monodromy coinvariants; primitivity therefore
identifies both kernels with the saturation of its image.
Consequently the exported matrices are matrices of the actual
collapse in suitable integral target bases.  This last argument also
eliminates any ambiguity from choosing splittings of the filtration.

\end{enumerate}
\end{proof}
